\section{Experimental Setup}
\label{sec:experimental_setup}

\subsection{Machine and Material}
\label{sec:machine}

All experiments use a Bantam Tools Desktop CNC Milling Machine, a 3-axis vertical milling center with a $140 \times 114 \times 38$~mm work envelope, 26,000~RPM maximum spindle speed, and 0.003~mm positioning resolution.
The workpiece material is ultra-high molecular weight polyethylene (UHMW-PE), selected for machinability and safe data collection (no hazardous chips or coolant).
All cuts use a 1/8\textquotesingle\textquotesingle\ (3.175~mm) 2-flute carbide end mill.

\subsection{Dataset}
\label{sec:dataset}

The dataset comprises 120 machining files spanning 9 operation classes: three toolpath types (face milling, pocket milling, adaptive clearing) crossed with three operating conditions.
Table~\ref{tab:dataset_classes} summarizes the class distribution.

\begin{table}[t]
\centering
\caption{Dataset composition: 9 operation classes, 120 files total. ``Standard'' denotes air-cut passes (spindle running, no workpiece engagement); ``150025'' denotes active-cut passes (0.025\textquotesingle\textquotesingle\ depth of cut in UHMW-PE); ``damage'' denotes air-cut with one spindle drive band deliberately removed to induce runout.}
\label{tab:dataset_classes}
\begin{tabular}{@{}llcc@{}}
\toprule
Toolpath & Condition & Files & Description \\
\midrule
Face milling    & Standard          & 20 & Flat surface generation \\
Face milling    & Modified (150025) & 20 & Altered feed/depth parameters \\
Pocket milling  & Standard          & 20 & Enclosed rectangular pocket \\
Pocket milling  & Modified (150025) & 20 & Altered feed/depth parameters \\
Adaptive clearing & Standard        & 20 & Constant-engagement roughing \\
Adaptive clearing & Modified (150025) & 20 & Altered feed/depth parameters \\
Face milling    & Tool damage       & 5  & Drive band removed (spindle runout) \\
Pocket milling  & Tool damage       & 5  & Drive band removed (spindle runout) \\
Adaptive clearing & Tool damage     & 5  & Drive band removed (spindle runout) \\
\midrule
\multicolumn{2}{@{}l}{\textbf{Total}} & \textbf{120} & \\
\bottomrule
\end{tabular}
\end{table}

The three toolpath types produce distinct G-code structures (different command sequences, coordinate patterns, parameter distributions) and distinct \emph{quasi-static} sensor signatures (trajectory-dependent magnetometer and drive-current patterns, thermal drift), providing a diverse evaluation ground.
We emphasize the regime explicitly: at the effective 4~Hz sampling rate (256 samples per 64-second window), the instrumentation captures program structure and slow process state, \emph{not} tooth-passing vibration---at 26{,}000~RPM with a 2-flute tool the tooth-passing frequency is ${\sim}867$~Hz, two orders of magnitude above the ${\sim}2$~Hz Nyquist limit. The detector therefore keys on program-structure and quasi-static physics rather than cutting-force or vibration spectra, and references to ``physics'' throughout should be read in that sense.

\paragraph{Sensor instrumentation.}
Six Arduino Nano 33 BLE Sense Lite boards are attached at different machine locations (spindle housing, worktable, column, base).
Each board provides 9-DOF inertial measurement (accelerometer, gyroscope, magnetometer), environmental sensors (temperature, humidity, pressure), audio level (microphone RMS), and color/proximity sensing (APDS-9960).
After feature engineering (spectral decomposition, windowed statistics, cross-channel correlations), each board contributes 18--24 channels.
An additional module provides spindle and axis drive current readings (8 channels), yielding 110 total channels per time step in the default (f110) configuration.
Table~\ref{tab:sensor_cost} itemizes the hardware cost.

\begin{table}[t]
\centering
\caption{Sensor hardware cost breakdown. The total system cost of \$492 is negligible relative to CNC machine costs (\$50K--\$500K).}
\label{tab:sensor_cost}
\begin{tabular}{@{}lrc@{}}
\toprule
Component & Quantity & Unit / Total Cost \\
\midrule
Arduino Nano 33 BLE Sense Lite & 6 & \$33 / \$198 \\
MCC DAQ USB-1808X (current measurement) & 1 & \$169 \\
Raspberry Pi 4 (data aggregation) & 1 & \$75 \\
Cabling, enclosures, mounting & --- & \$50 \\
\midrule
\textbf{Total} & & \textbf{\$492} \\
\bottomrule
\end{tabular}
\end{table}

\paragraph{Preprocessing.}
Sensor signals are segmented into 64-second windows (approximately one toolpath pass) and aligned with G-code tokens via timestamp matching.
V7 windows contain 6 tokens (3 content plus BOS, EOS, padding); V8 windows contain up to 64 tokens spanning multiple G-code lines.
Sensor features are standardized per-channel using training-set statistics.

Three feature configurations are evaluated:
\begin{itemize}
    \item \textbf{f110 (full)}: All 110 sensor channels.
    \item \textbf{f98}: 98 channels after removing features with pairwise $|r| > 0.95$.
    \item \textbf{f56}: 56 channels retaining only accelerometer, gyroscope, audio, and electrical modalities.
\end{itemize}

\paragraph{Dataset size.}
After windowing, the dataset comprises approximately 545 sensor windows across all 120 files---modest by deep learning standards, motivating cross-validation with bootstrap confidence intervals rather than single-split point estimates.

\subsection{Attack Generation Protocol}
\label{sec:attack_generation}

Synthetic attacks are generated by modifying the aligned G-code tokens for each sensor window according to Table~\ref{tab:attack_taxonomy}.
For each normal test window in each cross-validation fold, we generate one attacked variant per attack type and perturbation magnitude, yielding paired (normal, attacked) windows with identical sensor signals but different claimed G-code.

Attacks are generated deterministically (seed = 42).
A2 coordinate perturbation is evaluated across all three axes (X, Y, Z) at 10 magnitudes, yielding 30 sub-experiments per fold.
A1, A3, and A4 each produce a single deterministic attack variant per window.

\subsection{Evaluation Metrics}
\label{sec:metrics}

\paragraph{Primary metric: AUROC.}
The Area Under the ROC curve measures the probability that a randomly chosen attacked window scores higher than a randomly chosen normal window.
AUROC is threshold-free, making it appropriate for comparing methods without committing to a decision threshold.

\paragraph{Secondary metrics.}
\begin{itemize}
    \item \textbf{Detection rate at FPR = 5\%}: True positive rate when the false positive rate is held at 5\%.
    \item \textbf{Cohen's $d$}: Standardized mean difference between normal and attacked score distributions, $d = (\mu_{\text{attack}} - \mu_{\text{normal}}) / s_{\text{pooled}}$.
    \item \textbf{Expected Calibration Error (ECE)}: Weighted average absolute difference between predicted confidence and empirical accuracy across 15 equal-width bins.
\end{itemize}

\subsection{Statistical Methods}
\label{sec:statistical_methods}

\paragraph{Cross-validation.}
All results use 5-fold file-level stratified cross-validation.
File-level splitting prevents information leakage from temporally correlated sensor signals within a single file.

\paragraph{Bootstrap confidence intervals.}
95\% confidence intervals are computed by bootstrap resampling of test windows ($B = 5000$ iterations). Because windows within a file are temporally correlated, these window-level intervals can understate uncertainty relative to the effective (file-level) sample size; we therefore also report per-fold AUROC ranges (Appendix~\ref{app:perfold}) as a conservative cross-check.

\paragraph{Multiple comparison correction.}
We apply Holm--Bonferroni correction~\cite{holm1979simple} to control the family-wise error rate at $\alpha = 0.05$.
Pairwise AUROC comparisons use the DeLong test~\cite{delong1988comparing}.

\paragraph{Power considerations.}
With approximately 110 test windows per fold, the minimum detectable AUROC difference is approximately 0.08 at 80\% power (two-sided $\alpha = 0.05$, DeLong test variance).
We report exact $p$-values and confidence intervals throughout.
